En les centrals nuclears es produeix electricitat a partir de la fissió de nuclis d'urani. Aquesta energia s'utilitza per a generar vapor d'aigua, que fa girar una turbina. L'urani és un element químic metàl·lic de símbol U i nombre atòmic 92. A la natura trobem diferents isòtops de l'urani, però els més comuns són l'urani 238 i l'urani 235.

\begin{apartat}{1,25}
Calculeu el defecte de massa i l'energia d'enllaç per nucleó per a l'urani 235, i introduïu el valor obtingut a la taula de sota. Expliqueu la relació entre l'energia d'enllaç per nucleó i l'estabilitat del nucli. A partir d'aquí, indiqueu quin dels nuclis de la taula és el més estable.
\end{apartat}

\begin{center}
\begin{tabular}{|l|c|c|c|c|}
\hline
\textit{Nucli} & sofre 34 & ferro 56 & radi 226 & urani 235 \\
\hline
\textit{Energia d'enllaç per nucleó (MeV)} & 8,58 & 8,79 & 7,66 & \\
\hline
\end{tabular}
\end{center}

\begin{apartat}{1,25}
En una reacció nuclear de fissió de l'urani 235, un neutró d'alta energia impacta en un nucli d'urani. Com a resultat, es formen dos nuclis més petits i tres neutrons. Si considerem que un dels nuclis que es formen és el bari 141, escriviu-ne la reacció nuclear completa. Per a cada nucli d'urani fissionat, s'alliberen 202,5~MeV. Calculeu quants grams d'urani 235 són necessaris per a produir l'energia necessària per a il·luminar un estadi esportiu durant un partit en què es consumeixen aproximadament $25\,000\ \mathrm{kW\,h}$.
\end{apartat}

\dades{%
  $c = 3,00\times 10^{8}\ \mathrm{m\,s^{-1}}$.\\
  $1\ \mathrm{eV} = 1,602\times 10^{-19}\ \mathrm{J}$.\\
  Masses nuclears (en kg):\\[2pt]
  \begin{tabular}{|c|c|c|}
  \hline
  \textit{Protó} & \textit{Neutró} & \textit{Nucli d'urani 235} \\
  \hline
  $1,672\,622\times 10^{-27}$ & $1,674\,927\times 10^{-27}$ & $3,902\,158\times 10^{-25}$ \\
  \hline
  \end{tabular}\\[4pt]
  Nombre atòmic de diversos elements químics:\\[2pt]
  \begin{tabular}{|c|c|c|c|c|c|}
  \hline
  \textit{Kr} & \textit{Rb} & \textit{Sr} & \textit{Ba} & \textit{La} & \textit{U} \\
  \hline
  $Z = 36$ & $Z = 37$ & $Z = 38$ & $Z = 56$ & $Z = 57$ & $Z = 92$ \\
  \hline
  \end{tabular}}
